% TOP_PROBLEM: {"id": 6200067, "title": "Boundaries of Groups and Kleinian Groups — Problem 67", "category_id": 6, "category": {"id": 6, "name": "geometry", "display_name": "Geometry", "description": "Euclidean and non-Euclidean geometry, geometric structures.", "slug": "geometry", "order_index": 6, "created_at": "2026-07-31T15:26:25.671Z"}, "status": "partially_solved"}
% FIRST_POSTED: null
% ENTRY_KIND: statement_only
% Imported from: batch11.tex; UnsolvedMath ID 6200067
\documentclass[11pt]{article}
\usepackage[margin=1in]{geometry}
\usepackage{amsmath,amssymb,mathtools}
\usepackage{fontspec,unicode-math}
\setmainfont{Latin Modern Roman}
\setsansfont{Latin Modern Sans}
\setmathfont{Latin Modern Math}
\usepackage{xurl,hyperref}
\hypersetup{colorlinks=true,linkcolor=blue,urlcolor=blue}
\newcommand{\sref}[2]{\hyperlink{source:#1:#2}{[S#2]}}
\newcommand{\cataloguescope}[1]{\paragraph{Mathematical question.}#1}
\begin{document}
% UnsolvedMath ID 6200067; source catalogue code AMR-061-0067
\setcounter{section}{6200066}
\section{Boundaries of Groups and Kleinian Groups — Problem 67}\label{6200067}
{\small\sffamily UnsolvedMath ID 6200067 · Geometry}
\subsection{Definitions and mathematical statement}

\paragraph{Shared notation.}$\mathbb N=\{1,2,\ldots\}$, and $\mathbb Z,\mathbb Q,\mathbb R,\mathbb C$ denote the integers, rationals, reals and complex numbers. Any other conventions are specified locally.


A finitely generated group $G$ is word-hyperbolic if a Cayley graph for a finite generating set, with its word metric, has uniformly thin geodesic triangles: for some $\delta\geq0$, each side lies in the $\delta$-neighborhood of the union of the other two sides. Its Gromov boundary $\partial_\infty G$ is the set of geodesic rays based at the identity, modulo finite Hausdorff distance, with the topology of fellow traveling for longer and longer initial segments. Left multiplication gives the action of $G$ on this boundary.
Let $H\leq G$ be a subgroup that is itself word-hyperbolic in its own word metric. No quasiconvexity assumption on $H$ in $G$ is made. A map $F:\partial_\infty H\to\partial_\infty G$ is $H$-equivariant if $F(h\xi)=hF(\xi)$ for all $h\in H$ and $\xi\in\partial_\infty H$, using inclusion to define the target action. A Cannon--Thurston map is the boundary restriction of a continuous extension of the inclusion $H\hookrightarrow G$ to the Gromov compactifications; such a restriction is equivariant. For finite $H$ its boundary is empty and the map question is vacuous. The source's historical formulation asks the Cannon--Thurston existence question for every infinite hyperbolic subgroup.
\paragraph{Statement recorded in the supplied source.}
For every word-hyperbolic subgroup $H$ of a word-hyperbolic group $G$, must there exist an $H$-equivariant continuous map $\partial_\infty H\to\partial_\infty G$? \sref{6200067}{2}
\subsection{Short English statement}
Does every hyperbolic subgroup of a hyperbolic group admit a continuous map between their boundaries that respects the subgroup action, including when the subgroup is not quasiconvex?
\subsection{Sources}
\begin{description}
\item[\hypertarget{source:6200067:1}{S1}] ULAM.AI and the UnsolvedMath Contributors, UnsolvedMath: A Curated Collection of Open Mathematics Problems, record 6200067 (AMR-061-0067). CC BY 4.0. \url{https://huggingface.co/datasets/ulamai/UnsolvedMath}.
\item[\hypertarget{source:6200067:2}{S2}] Ilya Kapovich and Nadia Benakli, \emph{Boundaries of Hyperbolic Groups}, arXiv:math/0202286v1 (27 February 2002), \S2, Definitions 2.1, 2.4, 2.7--2.13, pp. 3--7; \S12, Definition 12.5 and final question before \S13, pp. 36--37. \url{https://arxiv.org/abs/math/0202286v1}.
\end{description}
\label{end:6200067}
\end{document}
